\chapter{サイクル数に関する考察}\label{chap:6}

\section{処理区間別のサイクル数}\label{sec:6-1}

固定乱数64要素では、8,614から5,527サイクルに減った（35.8\%減）。
READYまでの398サイクルは両方式で同じだった。差がついたのは主にREADYから両側のソートが終わるまでで、この区間は7,178から4,087サイクルになった（表\ref{tab:6}）。



\begin{table}[H]

\centering

\caption{乱数64要素のイベント間サイクル数}

\label{tab:6}

\footnotesize

\setlength{\tabcolsep}{3pt}

\renewcommand{\arraystretch}{1.25}

\begin{tabular}{>{\raggedright\arraybackslash}p{\dimexpr0.566821\linewidth-2\tabcolsep\relax}>{\raggedright\arraybackslash}p{\dimexpr0.216313\linewidth-2\tabcolsep\relax}>{\raggedright\arraybackslash}p{\dimexpr0.216865\linewidth-2\tabcolsep\relax}}\toprule

\textbf{観測区間} & \textbf{serial\_split} & \textbf{dual}\\\midrule

開始からREADY & 398 & 398\\

READYから両側ソート完了 & 7,178 & 4,087\\

両側完了からCOMPLETE & 1,037 & 1,039\\

COMPLETEから参加コア終端 & 1 & 3\\

総サイクル数 & 8,614 & 5,527\\

\bottomrule\end{tabular}

\end{table}

両側のソート完了からCOMPLETEまでは、1コアで1,037、2コアで1,039サイクルとほぼ同じだった。この区間にはマージと書き戻しに加え、フラグの確認と通知も含まれる。
初期コピーやマージはCore0が担当し、並行ソート中にも共有ポートの待ちがあるため、2コアでも総サイクル数は1コアの半分にはならない。



\section{入力配列・交換回数と総サイクル数}\label{sec:6-2}

\begin{table}[H]

\centering

\caption{64要素の2コア実行を入力別に比較}

\label{tab:7}

\footnotesize

\setlength{\tabcolsep}{3pt}

\renewcommand{\arraystretch}{1.25}

\begin{tabular}{>{\raggedright\arraybackslash}p{\dimexpr0.158229\linewidth-2\tabcolsep\relax}>{\raggedright\arraybackslash}p{\dimexpr0.160232\linewidth-2\tabcolsep\relax}>{\raggedright\arraybackslash}p{\dimexpr0.218316\linewidth-2\tabcolsep\relax}>{\raggedright\arraybackslash}p{\dimexpr0.266386\linewidth-2\tabcolsep\relax}>{\raggedright\arraybackslash}p{\dimexpr0.196836\linewidth-2\tabcolsep\relax}}\toprule

\textbf{入力} & \textbf{総数} & \textbf{両側完了\newline の時点} & \textbf{両側完了から\newline COMPLETE} & \textbf{Core1\newline ストール}\\\midrule

整列 & 4,398 & 3,481 & 914 & 388\\

逆順 & 5,391 & 4,473 & 915 & 389\\

乱数 & 5,527 & 4,485 & 1,039 & 925\\

\bottomrule\end{tabular}

\end{table}

乱数64要素は逆順64要素より136サイクル長い。
左右のソートが終わったのは、乱数入力で4,485サイクル目、逆順入力で4,473サイクル目で、差は12サイクルだった。そこからCOMPLETEまでにそれぞれ1,039サイクルと915サイクルを要し、この区間だけで124サイクルの差がついた。
差の大半は後段に現れており、ストール数だけでは総数を説明できない。



マージでは、左右の値の並びで比較回数や残り要素の転送が変わる。ただし後段にはフラグの確認と通知も含まれるため、124サイクルすべてがマージによる差とは言えない。
Core1のストール中にもCore0は動けるので、ストール回数をそのまま総サイクル数の増分として扱わない。



整列済み入力でも比較走査を続ける実装のため、4,398サイクルを要した。
Core1による配列へのstoreは、整列済み入力で0回、逆順入力で992回、乱数入力で468回だった。storeが最も多かったのは逆順入力だが、総サイクル数が最も多かったのは乱数入力だった。
ソート中の交換回数に加え、待機と後段処理を含む総数を見る必要がある。


